\documentclass[11pt]{article}
\usepackage[margin=1in]{geometry}
\usepackage{amsmath,amssymb,amsthm}
\usepackage{graphicx}
\usepackage{hyperref}

\newtheorem*{theorem}{Counterexample}
\newtheorem*{lemma}{Lemma}
\newtheorem*{definition}{Definition}

\newcommand{\R}{\mathbb R}
\newcommand{\E}{\mathbb E}
\newcommand{\Prob}{\mathbb P}
\newcommand{\ga}{\gamma}
\newcommand{\ginf}{\gamma_\infty}
\newcommand{\cL}{\mathcal L}

\title{A \(c_0\) counterexample to the gamma-multiplier strengthening}
\author{Run fa\_banach\_001, lane 5}
\date{2026-06-14}

\begin{document}
\maketitle

\section*{Source problem}

The source paper is Jan van Neerven, \emph{\(\gamma\)-Radonifying operators
-- a survey}, arXiv:0911.3788 \cite{vanNeerven2009}.  The question appears
on PDF page 24, immediately after Theorem 5.2, the
\(\gamma\)-bounded multiplier theorem.

\begin{center}
\includegraphics[width=\linewidth]{figures/open_problem_crop.png}
\end{center}

Theorem 5.2 says that if \(M:A\to\mathcal L(E,F)\) is strongly measurable
with \(\gamma\)-bounded range, then
\[
  \widetilde M:\gamma(L^2(A;H),E)\longrightarrow
  \gamma_\infty(L^2(A;H),F)
\]
is bounded.  The paragraph after the proof asks whether \(\widetilde M\)
actually takes values in the radonifying class \(\gamma\), already for
\(A=(0,1)\) and \(H=\R\).  The displayed paragraph writes the codomain as
\(\gamma(L^2(A;H),E)\); in the theorem's notation the evident codomain is
\(\gamma(L^2(A;H),F)\).  The counterexample below works even in the
same-space case \(E=F=c_0\).

\section*{Definitions used}

Let \(H\) be a Hilbert space and \(E\) a Banach space.  An operator
\(T:H\to E\) is \(\gamma\)-summing if
\[
  \|T\|_{\gamma_\infty(H,E)}^2
  =
  \sup_{(h_k)}
  \E\Big\|\sum_{k=1}^N \gamma_k T h_k\Big\|_E^2
  <\infty,
\]
where the supremum is over all finite orthonormal systems in \(H\).  It is
\(\gamma\)-radonifying if the corresponding Gaussian series converges in
\(L^2(\Omega;E)\) for an orthonormal basis; equivalently
\(\gamma(H,E)\) is the closed radonifying subspace of
\(\gamma_\infty(H,E)\).

A family \(\mathcal T\subseteq\mathcal L(E,F)\) is \(\gamma\)-bounded if
there is \(C\) such that for all finite choices \(T_j\in\mathcal T\) and
\(x_j\in E\),
\[
  \Big(\E\Big\|\sum_j\gamma_jT_jx_j\Big\|_F^2\Big)^{1/2}
  \le C
  \Big(\E\Big\|\sum_j\gamma_jx_j\Big\|_E^2\Big)^{1/2}.
\]

\section*{Proof Intuition}

The source survey already recalls the Linde--Pietsch diagonal operator
\[
  D:\ell_2\to c_0,\qquad
  D(\alpha_n)=\big(\alpha_n/\sqrt{\log(n+1)}\big)_n,
\]
which is \(\gamma\)-summing but not \(\gamma\)-radonifying.  The point is to
realize exactly this diagonal operator as \(T_{M\phi}\) for a multiplier
whose range is still \(\gamma\)-bounded.  Rademacher functions provide the
orthonormal coordinates in \(L^2(0,1)\), while the slow weights
\((\log(n+1))^{-1/2}\) are just small enough for all finite Gaussian maxima
to be uniformly square-integrable and just large enough that the full
Gaussian coordinate sequence fails to tend to zero.

\begin{lemma}[A Gaussian maximal estimate]\label{lem:max}
Let \(a_n=(\log(n+1))^{-1/2}\).  There is an absolute constant \(C\) such
that, for every finite or countable centred real Gaussian family \((Y_n)\) satisfying
\(\operatorname{Var}(Y_n)\le \sigma^2\) for all \(n\),
\[
  \E\sup_n |a_nY_n|^2 \le C^2\sigma^2 .
\]
\end{lemma}

\begin{proof}
It suffices first to bound finite suprema; the countable case then follows
by monotone convergence.  By scaling assume \(\sigma=1\).  For \(u\ge 3\), the Gaussian tail estimate
\(\Prob\{|Y_n|>v\}\le 2e^{-v^2/2}\) and the union bound give
\[
  \Prob\{\sup_n |a_nY_n|>u\}
  \le
  2\sum_{n\ge1}\exp\{-u^2\log(n+1)/2\}
  =
  2\sum_{n\ge1}(n+1)^{-u^2/2}.
\]
Since \(u\ge3\),
\[
  \sum_{n\ge1}(n+1)^{-u^2/2}
  \le
  2^{-(u^2/2-2)}\sum_{n\ge1}(n+1)^{-2}.
\]
Therefore
\[
  \E\sup_n |a_nY_n|^2
  =
  \int_0^\infty 2u\,\Prob\{\sup_n |a_nY_n|>u\}\,du
  \le
  9+
  C_0\int_3^\infty u\,2^{-u^2/2}\,du
  <\infty,
\]
with an absolute constant independent of the finite family.
\end{proof}

\begin{theorem}
The multiplier theorem in \cite[Theorem 5.2]{vanNeerven2009} cannot in
general be strengthened from \(\gamma_\infty\)-valued to
\(\gamma\)-valued, even for \(A=(0,1)\), \(H=\R\), and \(E=F=c_0\).
\end{theorem}

\begin{proof}
Let \((r_n)_{n\ge1}\) be the Rademacher functions on \((0,1)\), viewed as an
orthonormal sequence in \(L^2(0,1)\), and put
\[
  a_n=\frac{1}{\sqrt{\log(n+1)}},\qquad
  g(t)=(a_nr_n(t))_{n\ge1}\in c_0 .
\]
The map \(g:(0,1)\to c_0\) is strongly measurable, since it is the uniform
limit of its finite-coordinate truncations.

Let \(e_1^*:c_0\to\R\) be the first coordinate functional and define
\[
  M(t)x=e_1^*(x)g(t),\qquad x\in c_0.
\]
Then \(M:(0,1)\to\mathcal L(c_0,c_0)\) is strongly measurable.  Its range is
\(\gamma\)-bounded.  Indeed, for finite \(t_j\in(0,1)\) and \(x_j\in c_0\),
set \(b_j=e_1^*(x_j)\).  For each coordinate \(n\), the Gaussian variable
\(\sum_j\gamma_j b_j r_n(t_j)\) has variance at most \(\sum_j |b_j|^2\).
Lemma \ref{lem:max} gives
\[
  \E\Big\|\sum_j\gamma_j M(t_j)x_j\Big\|_{c_0}^2
  =
  \E\sup_n
  \Big|a_n\sum_j\gamma_j b_jr_n(t_j)\Big|^2
  \le
  C^2\sum_j |b_j|^2 .
\]
Since \(e_1^*\) is contractive,
\[
  \sum_j |b_j|^2
  =
  \E\Big|\sum_j\gamma_j e_1^*(x_j)\Big|^2
  \le
  \E\Big\|\sum_j\gamma_jx_j\Big\|_{c_0}^2.
\]
Thus \(\{M(t):0<t<1\}\) is \(\gamma\)-bounded.

Now take the constant finite-rank simple function
\[
  \phi(t):\R\to c_0,\qquad \phi(t)\lambda=\lambda e_1 .
\]
The induced operator \(T_{M\phi}:L^2(0,1)\to c_0\) is
\[
  T_{M\phi}f
  =
  \int_0^1 f(t)g(t)\,dt
  =
  \big(a_n\langle f,r_n\rangle_{L^2}\big)_{n\ge1}.
\]
Bessel's inequality implies \((\langle f,r_n\rangle)_n\in\ell_2\), hence
the displayed sequence is indeed in \(c_0\).

This operator is \(\gamma\)-summing.  If \(h_1,\ldots,h_N\) is an
orthonormal system in \(L^2(0,1)\), then the \(n\)-th coordinate of
\(\sum_k\gamma_kT_{M\phi}h_k\) is
\[
  a_n\sum_k\gamma_k\langle h_k,r_n\rangle .
\]
By Bessel's inequality its variance is at most \(a_n^2\).  Lemma
\ref{lem:max} therefore bounds
\[
  \E\Big\|\sum_k\gamma_kT_{M\phi}h_k\Big\|_{c_0}^2
\]
uniformly over all finite orthonormal systems.  Hence
\(T_{M\phi}\in\gamma_\infty(L^2(0,1),c_0)\), as Theorem 5.2 predicts.

It remains to show that \(T_{M\phi}\notin\gamma(L^2(0,1),c_0)\).  Let
\(R=\overline{\operatorname{span}}\{r_n:n\ge1\}\subseteq L^2(0,1)\).  If
\(T_{M\phi}\) were \(\gamma\)-radonifying on \(L^2(0,1)\), then its
restriction to \(R\) would be \(\gamma\)-radonifying.  Under the
identification \(R\simeq\ell_2\) sending \(r_n\) to the \(n\)-th unit vector,
this restriction is the diagonal map
\[
  D:\ell_2\to c_0,\qquad D(\alpha_n)=(a_n\alpha_n)_n .
\]
Thus the Gaussian series
\[
  \sum_{n\ge1}\gamma_n a_ne_n
\]
would converge in \(c_0\), and in particular the random sequence
\((a_n\gamma_n)_n\) would belong to \(c_0\) almost surely.

This is false.  The events
\[
  A_n=\{|\gamma_n|>\sqrt{\log(n+1)}\}
\]
are independent, and the standard lower Gaussian tail bound
\(\Prob\{|\gamma|>x\}\ge c e^{-x^2/2}/x\) for \(x\ge1\) yields
\[
  \sum_{n\ge1}\Prob(A_n)
  \ge
  c\sum_{n\ge1}\frac{1}{\sqrt{(n+1)\log(n+1)}}=\infty .
\]
By Borel--Cantelli, \(A_n\) occurs infinitely often almost surely.  Hence
\(|a_n\gamma_n|>1\) infinitely often almost surely, so
\((a_n\gamma_n)_n\notin c_0\).  This contradiction proves that
\(T_{M\phi}\notin\gamma(L^2(0,1),c_0)\).
\end{proof}

\section*{Verification notes}

The example satisfies the hypotheses of the source multiplier theorem:
\((0,1)\) is sigma-finite, \(H=\R\), \(E=F=c_0\), \(M\) is strongly
measurable with \(\gamma\)-bounded range, and \(\phi\) is constant finite
rank simple.  The induced operator lands in \(\gamma_\infty\) but not in
\(\gamma\), so it gives a negative answer to the proposed strengthening.

No computational verification is needed.  The script
\texttt{code/crop\_open\_problem.py} only regenerates the source crop.  A
bounded novelty check was performed on 2026-06-14: the run indexes were
searched for arXiv:0911.3788 and core terms including
\(\gamma\)-multiplier, \(\gamma_\infty\), \(c_0\), and radonifying; the
source section around Theorem 5.2 was inspected; web searches for the exact
open-problem phrase and for gamma-multiplier/c0 variants found no direct
answer.  The source itself contains the Linde--Pietsch diagonal
\(\gamma_\infty\setminus\gamma\) example; the extra observation here is that
Rademacher functions turn it into a multiplier with \(\gamma\)-bounded range.

\section*{Limitations}

This packet gives a negative answer in the target space \(c_0\), so it is
consistent with the positive statement noted in the source for target spaces
not containing an isomorphic copy of \(c_0\).  The novelty confidence is
moderate until a specialist checks whether this multiplier realization is
already folklore or appears implicitly elsewhere.

\section*{Human-review recommendation}

Check the \(\gamma\)-boundedness estimate for the rank-one family
\(M(t)=g(t)\otimes e_1^*\), and check that restricting a
\(\gamma\)-radonifying operator to the Rademacher subspace preserves
\(\gamma\)-radonification.  These are the two substantive steps; the rest is
the standard Linde--Pietsch \(c_0\) diagonal obstruction.

\begin{thebibliography}{9}
\bibitem{vanNeerven2009}
J. van Neerven,
\emph{\(\gamma\)-Radonifying operators -- a survey},
arXiv:0911.3788, 2009; revised 2010.
\end{thebibliography}

\end{document}
